\documentclass[11pt,a4paper]{article}

% Template visual Matriz Aprovação. Compile com XeLaTeX ou LuaLaTeX.
\usepackage{fontspec}
\usepackage{polyglossia}
\setmainlanguage{portuguese}
\IfFontExistsTF{Aptos}{
  \setmainfont{Aptos}
  \setsansfont{Aptos}
}{
  \IfFontExistsTF{Calibri}{
    \setmainfont{Calibri}
    \setsansfont{Calibri}
  }{
    \IfFontExistsTF{TeX Gyre Heros}{
      \setmainfont{TeX Gyre Heros}
      \setsansfont{TeX Gyre Heros}
    }{
      \setmainfont{Latin Modern Sans}
      \setsansfont{Latin Modern Sans}
    }
  }
}
\usepackage[a4paper,margin=2cm,headheight=16pt]{geometry}
\usepackage{graphicx}
\usepackage{xcolor}
\usepackage{tikz}
\usepackage{tcolorbox}
\usepackage{enumitem}
\usepackage{booktabs}
\usepackage{tabularx}
\usepackage{amssymb}
\usepackage{fancyhdr}
\usepackage{titlesec}
\usepackage{hyperref}

\definecolor{matrizlime}{HTML}{CBFF4D}
\definecolor{matrizink}{HTML}{101313}
\definecolor{matrizmuted}{HTML}{596363}
\definecolor{matrizline}{HTML}{DDE3DF}
\definecolor{matrizsoft}{HTML}{F3F7EC}

\hypersetup{colorlinks=true,linkcolor=matrizink,urlcolor=matrizink}
\pagestyle{fancy}
\fancyhf{}
\lhead{\textsf{\small MATRIZ APROVAÇÃO}}
\rhead{\textsf{\small APOSTILA}}
\cfoot{\textsf{\small \thepage}}
\renewcommand{\headrulewidth}{0.4pt}
\renewcommand{\headrule}{\hbox to\headwidth{\color{matrizline}\leaders\hrule height \headrulewidth\hfill}}

\titleformat{\section}{\Large\bfseries\color{matrizink}}{\thesection}{0.6em}{}
\titleformat{\subsection}{\large\bfseries\color{matrizink}}{\thesubsection}{0.6em}{}
\setlist{nosep,leftmargin=*}

\newtcolorbox{foco}{colback=matrizsoft,colframe=matrizlime,boxrule=1pt,arc=3pt,left=10pt,right=10pt,top=8pt,bottom=8pt}
\newtcolorbox{lei}{colback=matrizink,colframe=matrizink,coltext=white,boxrule=0pt,arc=3pt,left=10pt,right=10pt,top=8pt,bottom=8pt}
\newtcolorbox{alerta}{colback=white,colframe=matrizline,boxrule=0.6pt,arc=3pt,left=10pt,right=10pt,top=8pt,bottom=8pt}

% Logo usada na capa. Caminho relativo à pasta onde o .tex é compilado
% (materiais/<disciplina>/). Ajuste se a estrutura mudar.
\newcommand{\logomatriz}{../../public/logo.png}

\begin{document}

\begin{titlepage}
  \thispagestyle{empty}
  \begin{center}
    \includegraphics[width=0.55\linewidth]{\logomatriz}\par
  \end{center}
  \vspace{2.2cm}
  {\sffamily\fontsize{28}{32}\selectfont\bfseries Apostila — Informática: Segurança da Informação\par}
  \vspace{0.4cm}
  {\sffamily\large Receita Federal, CGU, PRF, Polícia Federal, INSS, TRT, TCU/TCE, Câmara e Petrobras\par}
  \vfill
  \begin{foco}
    \textsf{\textbf{Foco da apostila}}\\[3pt]
    
  \end{foco}
  \vspace{0.7cm}
  {\sffamily\small Atualizado em 16 de agosto de 2026\quad ·\quad Cebraspe, FCC e bancas semelhantes\par}
  \vspace{1cm}
  {\sffamily\small matrizaprova.com\par}
\end{titlepage}

\tableofcontents
\newpage

% Conteúdo gerado entra aqui.
\section*{O que cai}

Segurança da Informação costuma ser cobrada por conceitos, situações práticas e diferenças entre ameaças. Os pontos mais importantes são:

\begin{enumerate}
  \item princípios da segurança: confidencialidade, integridade, disponibilidade,
\end{enumerate}

   autenticidade e não repúdio;

\begin{enumerate}
  \item malware e formas de infecção;
  \item phishing, engenharia social e spoofing;
  \item senhas, autenticação multifator e controle de acesso;
  \item backup, criptografia e certificados digitais;
  \item diferença entre firewall, antivírus e mecanismos de detecção.
\end{enumerate}

\textbf{Atenção:} segurança não é apenas instalar antivírus. Ela combina pessoas, processos, controles físicos e tecnologias.

\section*{Teoria essencial}

\subsection*{1. Princípios da Segurança da Informação}

\subsubsection*{1.1. Confidencialidade}

Confidencialidade é garantir que a informação seja acessada somente por pessoas ou sistemas autorizados. Controle de permissões, criptografia e classificação da informação são mecanismos relacionados a esse princípio.

Uma quebra de confidencialidade ocorre quando um arquivo sigiloso é copiado ou visualizado por pessoa sem autorização, ainda que o conteúdo não seja alterado.

\subsubsection*{1.2. Integridade}

Integridade é preservar a exatidão, a completude e a não alteração indevida da informação. Hashes, assinaturas digitais e controle de versões ajudam a verificar se um arquivo foi modificado.

Se alguém altera o valor de uma nota em um sistema sem autorização, há quebra de integridade. O arquivo pode continuar disponível e acessível, mas seu conteúdo deixou de ser confiável.

\subsubsection*{1.3. Disponibilidade}

Disponibilidade é permitir que usuários autorizados acessem a informação e os serviços quando necessários. Redundância, manutenção, nobreak, backups, monitoramento e proteção contra ataques de negação de serviço contribuem para esse princípio.

\subsubsection*{1.4. Autenticidade}

Autenticidade permite verificar a identidade de uma pessoa, sistema ou origem de uma mensagem. Login, certificado digital e assinatura digital podem ajudar nesse processo.

Autenticidade não é a mesma coisa que autorização. A primeira responde "quem é?"; a segunda responde "o que essa identidade pode fazer?".

\subsubsection*{1.5. Não repúdio}

Não repúdio busca impedir que o autor de uma ação negue posteriormente sua realização. Assinaturas digitais e registros confiáveis de auditoria podem contribuir para esse objetivo.

\subsection*{2. Malware}

Malware é software malicioso criado para causar dano, obter acesso indevido, roubar informações ou explorar recursos do dispositivo.

\begin{tabularx}{\linewidth}{lX}
\toprule
 \textbf{Tipo} & \textbf{Característica} \\
\midrule
 vírus & depende de arquivo hospedeiro e normalmente precisa de alguma ação para se propagar \\
 worm & propaga-se automaticamente por redes ou sistemas vulneráveis \\
 trojan & disfarça-se de programa legítimo para induzir a instalação \\
 ransomware & bloqueia ou cifra dados e exige pagamento ou outra vantagem \\
 spyware & monitora atividades e coleta informações \\
 keylogger & registra teclas digitadas, inclusive senhas \\
 adware & exibe publicidade indesejada; pode também rastrear hábitos \\
 rootkit & busca ocultar presença e manter acesso privilegiado \\
 bot & transforma o dispositivo em parte de uma rede controlada \\
\bottomrule
\end{tabularx}


Um arquivo infectado por vírus pode contaminar outros quando é executado ou compartilhado. O worm não precisa necessariamente de um arquivo hospedeiro e é conhecido pela capacidade de propagação automática.

\subsection*{3. Engenharia social e phishing}

Engenharia social explora comportamento humano para obter acesso, informação ou vantagem. O atacante pode usar autoridade falsa, urgência, medo, curiosidade ou recompensa.

\textbf{Phishing} é uma fraude que tenta induzir a vítima a clicar em link, abrir arquivo ou fornecer dados, geralmente por e-mail, mensagem ou página falsa.

Variações comuns:

\begin{itemize}
  \item \textbf{spear phishing:} ataque direcionado a uma pessoa ou organização;
  \item \textbf{smishing:} phishing por SMS ou aplicativo de mensagens;
  \item \textbf{vishing:} fraude por ligação ou voz;
  \item \textbf{pharming:} redirecionamento para página falsa, mesmo quando a vítima tenta acessar um endereço legítimo.
\end{itemize}

\textbf{Spoofing} é falsificação de identidade, origem ou endereço. Pode envolver e-mail, número de telefone, endereço IP, DNS ou página visualmente semelhante.

Cuidados: conferir domínio, não confiar apenas no nome exibido, evitar links inesperados, confirmar solicitações por canal independente e nunca fornecer senha ou código de autenticação por mensagem.

\subsection*{4. Senhas e autenticação}

Autenticação é o processo de verificar a identidade. Os fatores de autenticação são normalmente classificados como:

\begin{itemize}
  \item algo que o usuário \textbf{sabe}: senha ou PIN;
  \item algo que o usuário \textbf{possui}: token, celular ou cartão;
  \item algo que o usuário \textbf{é}: impressão digital ou reconhecimento facial.
\end{itemize}

Autenticação multifator combina fatores diferentes. Digitar senha e depois confirmar um código em aplicativo usa conhecimento e posse. Usar duas senhas continua sendo, em regra, um único fator: algo que o usuário sabe.

Uma boa senha deve ser exclusiva, longa e não baseada em informações públicas. O armazenamento seguro deve usar funções de derivação e mecanismos adequados, nunca guardar senhas em texto puro.

\subsection*{5. Criptografia, hash e assinatura digital}

Criptografia transforma uma informação legível em conteúdo protegido por meio de algoritmo e chave.

\subsubsection*{5.1. Criptografia simétrica}

Usa a mesma chave para cifrar e decifrar. É eficiente para grandes volumes de dados, mas exige que a chave seja compartilhada com segurança.

\subsubsection*{5.2. Criptografia assimétrica}

Usa um par de chaves relacionadas: pública e privada. A chave pública pode ser compartilhada; a privada deve permanecer protegida.

Para confidencialidade, uma mensagem pode ser cifrada com a chave pública do destinatário, que a decifra com sua chave privada. Para assinatura digital, usa-se a chave privada do signatário, e a verificação é feita com a chave pública correspondente.

\subsubsection*{5.3. Hash}

Hash é uma função que produz um resumo de tamanho definido a partir de uma entrada. Uma alteração no arquivo tende a produzir outro resumo. Hash não é criptografia reversível: não existe, em condições normais, uma chave para "desfazer" o hash.

\subsubsection*{5.4. Assinatura digital}

Assinatura digital ajuda a verificar autoria, integridade e, em determinados contextos, não repúdio. Ela não é simplesmente uma imagem da assinatura manuscrita colada em um documento.

\subsection*{6. Firewall, antivírus e IDS/IPS}

\begin{itemize}
  \item \textbf{Firewall:} controla tráfego de rede com base em regras. Pode bloquear ou permitir conexões de acordo com origem, destino, porta, protocolo e outros critérios.
  \item \textbf{Antivírus/antimalware:} identifica, bloqueia ou remove códigos maliciosos.
  \item \textbf{IDS:} detecta e alerta sobre atividades suspeitas.
  \item \textbf{IPS:} além de detectar, pode bloquear ou impedir a atividade identificada.
\end{itemize}

Nenhum desses mecanismos garante proteção absoluta. Segurança depende de configuração, atualização, monitoramento e comportamento seguro.

\subsection*{7. Backup}

Backup é uma cópia de dados para recuperação em caso de perda, corrupção, falha, ataque ou exclusão acidental.

A estratégia \textbf{3-2-1} recomenda manter:

\begin{itemize}
  \item pelo menos 3 cópias dos dados;
  \item em 2 tipos de mídia diferentes;
  \item pelo menos 1 cópia fora do ambiente principal.
\end{itemize}

Backup conectado permanentemente ao computador pode ser atingido por ransomware. É necessário testar a restauração: uma cópia que nunca foi recuperada pode não ser utilizável quando necessária.

\section*{Como a banca cobra}

\begin{itemize}
  \item Confidencialidade trata de acesso; integridade trata de alteração; disponibilidade trata de acesso no momento necessário.
  \item Autenticação identifica; autorização define permissões.
  \item Hash não é criptografia reversível.
  \item MFA exige fatores diferentes, não apenas duas senhas.
  \item Worm se propaga automaticamente; trojan depende de disfarce e indução.
  \item Firewall controla tráfego; antivírus atua contra código malicioso.
  \item Backup é cópia de recuperação, não simplesmente sincronização.
\end{itemize}

\section*{Exemplos resolvidos}

\subsection*{Exemplo 1}

Um atacante modifica o valor de uma transferência bancária, mas o usuário continua conseguindo acessar o sistema. Qual princípio foi violado?

\textbf{Integridade.} O dado foi alterado indevidamente. A disponibilidade foi mantida, e não há informação suficiente para afirmar quebra de confidencialidade.

\subsection*{Exemplo 2}

Uma mensagem solicita que o usuário clique em um link para evitar o bloqueio imediato da conta. A página imita o banco e pede senha e código. Qual é o ataque?

\textbf{Phishing}, com uso de engenharia social. A urgência é usada para induzir a vítima a fornecer credenciais em uma página falsa.

\section*{Quadro-resumo}

\begin{tabularx}{\linewidth}{lX}
\toprule
 \textbf{Conceito} & \textbf{Pergunta principal} \\
\midrule
 confidencialidade & quem pode acessar? \\
 integridade & o conteúdo foi alterado? \\
 disponibilidade & o serviço está acessível quando necessário? \\
 autenticidade & quem é o autor ou usuário? \\
 autorização & o que ele pode fazer? \\
 vírus & usa arquivo hospedeiro \\
 worm & propaga-se automaticamente \\
 trojan & disfarça-se de programa legítimo \\
 ransomware & bloqueia ou cifra dados \\
 hash & verifica resumo, não é reversível \\
\bottomrule
\end{tabularx}


\section*{Questões de fixação}

\subsection*{Questão 1 — (questão inédita)}

Um usuário autorizado acessa um arquivo, mas encontra seu conteúdo alterado sem permissão. O principal princípio afetado é:

A) confidencialidade. \\ B) integridade. \\ C) disponibilidade. \\ D) autenticidade. \\ E) não repúdio.

\begin{alerta}
\textbf{Gabarito: B.} Integridade protege a exatidão e a não alteração indevida da informação.
\end{alerta}

\subsection*{Questão 2 — (questão inédita)}

A combinação de senha e código gerado por aplicativo caracteriza:

A) duas senhas do mesmo fator. \\ B) autenticação multifator. \\ C) criptografia simétrica. \\ D) assinatura digital. \\ E) autorização sem autenticação.

\begin{alerta}
\textbf{Gabarito: B.} A senha é algo que o usuário sabe; o aplicativo ou dispositivo é algo que possui. São fatores diferentes.
\end{alerta}

\subsection*{Questão 3 — (questão inédita)}

Assinale a alternativa correta sobre malware.

A) Worm sempre precisa de arquivo hospedeiro. \\ B) Trojan é um mecanismo de backup. \\ C) Ransomware pode cifrar dados e exigir pagamento. \\ D) Spyware é um protocolo de criptografia. \\ E) Rootkit é um equipamento de rede.

\begin{alerta}
\textbf{Gabarito: C.} Ransomware pode impedir o acesso aos dados por meio de criptografia e exigir resgate.
\end{alerta}

\subsection*{Questão 4 — (questão inédita)}

Sobre hash, é correto afirmar que:

A) sempre permite recuperar o arquivo original com a chave pública. \\ B) é sinônimo de criptografia simétrica. \\ C) gera um resumo usado, entre outras finalidades, para verificar alterações. \\ D) impede qualquer tipo de malware. \\ E) exige duas senhas para funcionar.

\begin{alerta}
\textbf{Gabarito: C.} Hash produz um resumo e pode ser usado para verificar integridade; não é uma cifra reversível.
\end{alerta}

\subsection*{Questão 5 — (questão inédita)}

Na estratégia de backup 3-2-1, recomenda-se:

A) 3 cópias, em 2 tipos de mídia, com 1 fora do ambiente principal. \\ B) 3 senhas, em 2 usuários, com 1 administrador. \\ C) 3 firewalls, em 2 redes, com 1 cópia local. \\ D) 3 antivírus, em 2 computadores, com 1 arquivo. \\ E) 3 cópias exclusivamente online.

\begin{alerta}
\textbf{Gabarito: A.} A estratégia combina redundância, mídias diferentes e uma cópia fora do ambiente principal.
\end{alerta}

\section*{Revisão final}

\begin{itemize}
  \item[$\square$] Diferencio confidencialidade, integridade e disponibilidade.
  \item[$\square$] Sei separar autenticação de autorização.
  \item[$\square$] Reconheço vírus, worm, trojan, ransomware e spyware.
  \item[$\square$] Identifico phishing, smishing, vishing e spoofing.
  \item[$\square$] Sei o que caracteriza autenticação multifator.
  \item[$\square$] Diferencio criptografia simétrica e assimétrica.
  \item[$\square$] Sei que hash não é reversível como uma cifra.
  \item[$\square$] Diferencio firewall, antivírus, IDS e IPS.
  \item[$\square$] Sei aplicar a regra de backup 3-2-1.
\end{itemize}

\end{document}
